\section{Scope, analytic premises, and what has been established}
\label{sec:premises}

The principal point-separation assertion A1 is proved here without
number-field GRH, by a complete checked reconstruction of the later
ANABELIAN\_A1 argument. The originally proposed passage from Hooley's
primitive-root density to theta-divisor incidence is not a theorem:
A04 identifies the logical gap, and A06 gives an explicit geometric
counterexample to its universal primitive-root version. This changes
the starting diagnosis. It does not turn the counterexample into a
counterexample to A1.

The return is \textbf{INCOMPLETE as an exhaustive audit}; the
requested proof through Hooley is also not established. A1 itself
is proved by the different route just specified. In source coverage,
800 of the 1,190 full register statements were recovered, and 591
statement/proof blocks are available. The 210-row
ANABELIAN\_A1\_SUCCESSOR and 180-row GT\_TOWER\_CHECK registers were not
recovered; their visible excerpts are audited separately. Availability
of a proof is not a claim that its whole content was checked. Each
actual check is recorded in \texttt{CHECK\_EVENTS.tsv}. The positive-rank
theta-incidence variant and unrelated open hypotheses in programme
rows remain recorded as open, rather than being assigned consequences
of GRH.

\subsection{The two inputs are deliberately different}

\paragraph{Input A: the stipulated Alp\"oge theorem.}
We accept the cited Dirichlet GRH theorem as instructed. The exact
supplied declarations in \texttt{ALPOGE\_Main.lean} are:
\begin{quote}\small\raggedright
\texttt{riemannHypothesis}\\
\texttt{riemannZeta\_eq\_zero\_re}\\
\texttt{LFunction\_ne\_zero\_of\_half\_lt\_re}\\
\texttt{LFunction\_eq\_zero\_re}\\
\texttt{LFunction\_eq\_zero\_re\_of\_isPrimitive}.
\end{quote}
The definitions in \texttt{ALPOGE\_Statements.lean} quantify over
Dirichlet characters of modulus $N$; the critical-strip assertion
has $0<\Re(s)<1$ and concludes $\Re(s)=1/2$.
The note's summary GRH claim is step~20, p.~4 \cite{Alpoge}.
No other intermediate step of the four-page note is needed in the
proofs below. In particular we do not convert its phrases about
purity into new geometric theorems.

This run inspected those declarations, definitions, and the attached
scan report. It did not build the Lean archive or rescan 175,187
source files. Acceptance as a cited theorem is the convention of
this task, not a newly reported verification of its formal proof.

\paragraph{Input M: the Masullo draft claim.}
Write M, or $\mathrm H_M$, for the draft's claimed GRH for the Hecke
$L$-functions of every number field, as described in \S1.2
\cite{Masullo}. Every use is explicitly conditional on that claim.
The trivial Hecke character over $F$ gives $\zeta_F$, so M includes
Dedekind GRH for every field used below.

\paragraph{Finding S01 (the exact numbered-source boundary).}
The supplied draft does not contain a completed numbered GRH
statement at its intended culminating locator. On printed/PDF
page~522, Chapter~13, Theorem~13.3.1.1 is a trace formula, referring
back to Theorem~12.5.1.1. The next heading is
Theorem~13.3.1.2, \emph{Divisor of $L(s,\chi)$}, but it has no
statement. Its preceding cross-reference is unresolved, and the
footnotes ask for insertion of the statement and proof of the main
theorems. Therefore this report cites M as the claim of \S1.2, with
13.3.1.2 as its unfinished intended location; it never cites the
latter as an already stated and proved GRH theorem.

\begin{proof}
The page was inspected both by direct PDF text extraction and by
rendering the page. The displayed material under 13.3.1.1 is the
trace formula. After the heading 13.3.1.2 the page has only the
editorial notes, and Chapter~14 starts on the following page.
This proves the source-status assertion. It neither supplies a
mathematical refutation of M nor claims a review of all 940 pages.
\end{proof}

\paragraph{Finding S02 (the analytic dependency rule).}
Input A discharges Dedekind GRH for finite abelian extensions of
$\mathbb Q$. M discharges it for arbitrary number fields. Neither
input, as used here, supplies a missing geometric existence theorem,
a missing finite-ramification hypothesis, or a description of an
unspecified finite quotient. No Artin holomorphy conjecture is needed
for the Chebotarev bounds chosen in this report.

\begin{proof}
For an abelian extension $F/\mathbb Q$ the Dedekind zeta function
is the product of the primitive Dirichlet $L$-functions corresponding to the
characters of $\operatorname{Gal}(F/\mathbb Q)$. At an unramified
prime with Frobenius order $f$, the character product of Euler
factors is $(1-p^{-fs})^{-[F:\mathbb Q]/f}$, exactly the Dedekind
factor; the inertia-invariant character factors give the ramified
version. Their nontrivial
zeros therefore satisfy the asserted line condition under A.
For arbitrary $F$, the trivial character specialization of M is
exactly $\zeta_F$. The remaining sentences specify hypotheses not
included in either premise; the concrete failures to make the
corresponding inferences are proved in A04, A06, T02 and G04.
The exact Chebotarev theorem G-CHEB assumes only the indicated
Dedekind zeta RH statements.
\end{proof}

There is a necessary fourth status in the audit: \emph{already
unconditional}. A qualitative Chebotarev theorem, a finite-field
Weil bound, an elementary modular computation, or a previously proved
geometric result does not become unconditional for the first time
because A is added. Calling such a result ``unconditional by Alp\"oge''
would misstate its dependency. The row tables retain the requested
A-only/M-only/beyond-both distinctions for actual analytic upgrades,
and this fourth status for existing statements.

\subsection{The geometric framework and exact locators}

All finite number-field monodromy quotients below are finite etale
covers of an open subscheme of $\operatorname{Spec}\mathbb Z$ (or
$\operatorname{Spec}\mathcal O_K$). The cyclotomic examples are the
covers obtained by adjoining roots of unity. The underlying
fibre-functor and generic-point correspondences are SGA~1,
Expos\'e~V, Th\'eor\`eme~4.1, \S7 and Propositions~8.1--8.2.
For a geometrically connected curve, its arithmetic fundamental-group
exact sequence and point sections are SGA~1, IX, Th\'eor\`eme~6.1
and Corollaire~6.4. Purely inseparable invariance is IX,
Th\'eor\`eme~4.10. These are used with the arithmetic Galois quotient
unchanged.

The configuration argument in A05 uses smooth proper curve lifting
(SGA~1, III, Th\'eor\`eme~7.3 and Corollaire~7.4), with marked
points then lifted by smoothness over the complete henselian base,
proper smooth
specialization (X, Th\'eor\`eme~3.8 and Corollaire~3.9), and the
prime-to-characteristic comparison and inertia of smooth proper pairs
(XIII, Corollaire~2.9, \S2.10, Lemme~2.11). Its surface relation is
XIII, Corollaire~2.12, and its tame Kummer extension at a boundary
uses XIII, Proposition~5.2 \cite{SGA1}. The exact archive context is
cote~149, scans~33--36 for Kummer separation and its covering
criterion, and scans~37--38 for the surrounding geometric reductions.

For compactified curves over finite fields, A05 uses the trace
formula and purity of $H^1$, with Frobenius eigenvalues of weight one:
Deligne, \emph{La conjecture de Weil. I}, formula~(1.5.1) and
Th\'eor\`eme~1.6 \cite{WeilI}. The equivalent action on the
Tate module and its dual must retain the full integral adic module,
not just its reduction modulo a prime. The curve cohomology and
orientation used there are also explicitly given in
\cite[Proposition~14.2 and Remark~14.10]{MilneLEC}.

The coherent Frobenius calculation in A03 and A06 concerns the
Artin--Schreier sequence (\cite[Example~7.9(b)]{MilneLEC}) and the stable part of $H^1(C,\mathcal O_C)$.
It is a characteristic-$p$ calculation of $p$-rank. No theorem about
complex zeros of number-field $L$-functions changes that rank.
Finite quotient representations in G01--G03 are weight-zero Artin
realizations; the original Tate modules may have positive weight.
No new purity theorem, motivic standard conjecture, crystalline
comparison, divided-power identity, or $\lambda$-ring theorem is
asserted. Those additional structures are not inputs to the proofs
in this report.
